% No preamble. Prefix dgb-. Exact finite identities, not a uniform comparison.
\section{An exact gauge and coarea bridge for the actual determinant}
\label{dgb-section}

This section rewrites the actual leading correction as coupled
vector, ghost and coarse-constraint determinants. It keeps the
constraint-current curvature and every finite orbit-density factor.
It also exhibits a soft mode that prevents a separate uniformly
massive estimate for the based ghost. The bounded finite-part
comparison with a standard background-Feynman determinant remains
a separate estimate.

\subsection{The actual multiplier Hessian}
\label{dgb-setup-section}

Fix a finite compatible regulator and endpoint $K=\ell^k$.
Let $U=V_*^{M,K}(W)$ be the physical stationary representative,
so $F_K(U)=W$ and $P_K(U)=I$. All matrices use the actual
current-link angle and scalar-site counting coordinates in a real
trace-orthonormal color basis.

In left-relative coordinates define
\begin{equation}
 s(x)=S_M(e^{ix}U),\qquad
 f(x)=\frac1i\log\!\left[F_K(e^{ix}U)W^{-1}\right],
 \qquad C=f'(0).
 \label{dgb-relative-coordinates}
\end{equation}
The target coordinate is a relative increment at $W$, not its
ordinary logarithm. Constrained stationarity gives a unique
multiplier $\lambda$ with $s'(0)=C^T\lambda$. Put
\begin{equation}
 \mathcal L=s''(0)-\sum_c\lambda_c f_c''(0).
 \label{dgb-Lagrangian}
\end{equation}
This is the Hessian of the critical Lagrangian
$s-\lambda\cdot f$. Its second term is the full
constraint-current curvature.

Let $C_p$ be the invertible actual pivot block, $I_p$ the pivot
insertion, and $E$ the retained insertion. Write
\begin{equation}
 h=I_pC_p^{-1},\qquad K_0=E-hCE,\qquad
 Ch=I,\qquad CK_0=0.
 \label{dgb-chart-linear}
\end{equation}
The subscript distinguishes $K_0$ from the endpoint ratio $K$.
The nonlinear source chart satisfies
\begin{equation}
 \Phi'(0)=K_0,\qquad
 \Phi''(0)[u,v]=-h f''(0)[K_0u,K_0v].
 \label{dgb-chart-curvature}
\end{equation}
Consequently
\[
 D_z^2s(\Phi(z))\big|_{z=0}=K_0^T\mathcal L K_0 .
\]
The actual absolute contour penalty vanishes at the saddle;
its Hessian is $T^TT$, where $T$ differentiates the contour
logarithms in the established normalization. Thus
\begin{equation}
 A_T=K_0^T(\mathcal L+T^TT)K_0>0.
 \label{dgb-actual-precision}
\end{equation}
Apart from field-independent constants, the actual leading
amplitude is
\begin{equation}
                  |\det C_p|^{-1}\det(A_T)^{-1/2}.
 \label{dgb-actual-amplitude}
\end{equation}
The real determinant signs are fixed from the flat convention.
All logarithms below are normalized at $W=1$ and then continued
on their chosen holomorphic branches.

\subsection{The based gauge nullspace}
\label{dgb-nullspace-section}

Let $D_b$ send scalar site parameters vanishing at every
$K$-center to their left-relative infinitesimal gauge variations
at $U$. Exact based invariance of the ordinary constraint gives
\begin{equation}
                              CD_b=0.
 \label{dgb-based-constraint}
\end{equation}
The scalar $s-\lambda\cdot f$ is based-gauge invariant and is
critical at zero. Differentiating its gauge identity there gives
\begin{equation}
                              \mathcal L D_b=0.
 \label{dgb-based-null}
\end{equation}
This need not hold for $s''$ without its multiplier-curvature
term. The actual contour endpoint covariance gives
\begin{equation}
                              TD_b=-I
 \label{dgb-contour-gauge}
\end{equation}
with the standard orientation. Changing the generator sign
changes only this sign. A fixed color normalization in the
gauge penalty contributes only a background-independent
factor, canceled by normalization.

The based action is free, hence $D_b$ is injective and
\begin{equation}
                              G_b=D_b^TD_b>0.
 \label{dgb-based-ghost}
\end{equation}
On $\ker C$, the form $\mathcal L$ is nonnegative, has
nullspace precisely $\operatorname{ran}D_b$, and is positive
on its quotient. Indeed, for any $x\in\ker C$, use
\eqref{dgb-contour-gauge} to add a unique based gauge vector
making $Tx=0$. Equation \eqref{dgb-based-null} leaves its
$\mathcal L$ value unchanged. On this contour-zero complement,
the positive form \eqref{dgb-actual-precision} equals
$\mathcal L$. Equality occurs only for a vertical vector.

\subsection{A finite Gaussian gauge-change identity}
\label{dgb-gauge-change-section}

\begin{lemma}[Gauge replacement with its orbit density]
\label{dgb-gauge-lemma}
Let $B_0$ be a symmetric nonnegative form on a finite space
$Z$, with nullspace $\operatorname{ran}V$, positive on the
quotient. Suppose $F_1,F_2:Z\to\mathbb R^s$ have invertible
$F_jV$. Then
\begin{equation}
 |\det(F_1V)|\det(B_0+F_1^TF_1)^{-1/2}
   =|\det(F_2V)|\det(B_0+F_2^TF_2)^{-1/2}.
 \label{dgb-gauge-identity}
\end{equation}
\end{lemma}

\begin{proof}
Choose one complement $E_0$ to $\operatorname{ran}V$ and set
$z=E_0u+V\alpha$. The change-of-variable Jacobian is the same
in the two integrals. The $B_0$ term depends only on $u$.
Integration over $\alpha$ yields
$(2\pi)^{s/2}/|\det(F_jV)|$, independently of $u$, by
translation and an invertible linear change of variables.
The remaining integral in $u$ is common. This proves
\eqref{dgb-gauge-identity}; no orbit-coordinate density
has been discarded.
\end{proof}

Apply the lemma on retained coordinates, with vertical insertion
$V$ characterized by $K_0V=D_b$. Take
$F_1=TK_0$ and $F_2=D_b^TK_0$. Equations
\eqref{dgb-contour-gauge} and \eqref{dgb-based-ghost} give
$F_1V=-I$ and $F_2V=G_b$. Therefore
\begin{equation}
 \det(A_T)^{-1/2}
     =\det(G_b)\det(A_b)^{-1/2},\qquad
 A_b=K_0^T(\mathcal L+D_bD_b^T)K_0.
 \label{dgb-divergence-replacement}
\end{equation}
This replaces the actual contour gauge by based covariant
divergence. It does not yet replace it by a full-site massive
background-Feynman gauge.

\subsection{Coarea removes the pivot}
\label{dgb-coarea-section}

Choose any symmetric coarse matrix $B$ such that
\begin{equation}
 Q_b=\mathcal L+D_bD_b^T+C^TBC>0,\qquad
 M_b=CQ_b^{-1}C^T>0.
 \label{dgb-completion}
\end{equation}
A completion exists at each fixed finite saddle: the first
two terms are positive on $\ker C$, so a sufficiently large
scalar $B$ makes the complementary Schur block positive.
The condition is positivity of $Q_b$. Merely requiring
$B>0$ need not suffice; conversely $B$ itself may be indefinite.

Let $n$ and $m$ be the ambient and constraint dimensions.
The same constrained Gaussian has the exact evaluations
\begin{align}
 \int_{\mathbb R^n}\delta(Cx)e^{-x^TQ_bx/2}\,dx
  &=(2\pi)^{(n-m)/2}|\det C_p|^{-1}\det(A_b)^{-1/2}
       \notag\\
  &=(2\pi)^{(n-m)/2}
           [\det Q_b\,\det M_b]^{-1/2}.
 \label{dgb-coarea-integral}
\end{align}
For the first, write $x=K_0z+hy$. In retained/pivot order,
this transformation has determinant $(\det C_p)^{-1}$,
and $Cx=y$. For the second, the normalized full Gaussian
of covariance $Q_b^{-1}$ sends $Cx$ to a Gaussian of covariance
$M_b$; its density at zero gives the factor.
Equivalently,
\begin{equation}
 \det Q_b\,\det(CQ_b^{-1}C^T)
                         =(\det C_p)^2\det A_b .
 \label{dgb-coarea-determinant}
\end{equation}

The delta in \eqref{dgb-coarea-integral} uses the target's
relative Lie coordinate. The actual target measure is group
Haar, whose density at that coordinate origin is the fixed
$h_0$. It supplies no factor $j(\log W)$.
If one instead uses induced Euclidean measure on $\ker C$,
its coarea factor is precisely included in
\eqref{dgb-coarea-integral}.

For a scalar expression $F$, write
$\mathcal N_0F=F(W)-F(1)$. Combining
\eqref{dgb-actual-amplitude},
\eqref{dgb-divergence-replacement} and
\eqref{dgb-coarea-integral} proves the exact identity
\begin{equation}
 \boxed{\displaystyle
 D_{M,K}
   =\mathcal N_0\left[
       \log\det G_b-\frac12\log\det Q_b
                    -\frac12\log\det M_b\right].}
 \label{dgb-main-identity}
\end{equation}
The pivot has disappeared, but the coarse-constraint
determinant remains. Every completion allowed in
\eqref{dgb-completion} gives the same combined expression.
The individual terms need not be completion-independent
or locally gauge invariant.

At fixed regulator these are identities on a positive real
germ and then on their chosen analytic branches. The algebraic
identity \eqref{dgb-coarea-determinant} also continues wherever
its inverses exist. No uniform complex domain for the three
individual logarithms, as $K$ or volume grows, follows from
the combined identity.

\subsection{A positive completion from the actual minimum action}
\label{dgb-residual-section}

In the same relative target coordinate set
\[
 s_{\rm eff}(y)=S_{MK}(e^{iy}W),\qquad
 q_{\rm eff}=s_{\rm eff}''(0).
\]
The envelope identity gives $s_{\rm eff}'(0)=\lambda$.
The residual scalar
\[
                 \mathcal R(x)=s(x)-s_{\rm eff}(f(x))
\]
is critical and nonnegative near zero by the actual local
minimum comparison. It is invariant under the full local
fine gauge group, because both scalar actions are locally
gauge invariant and the composite map is equivariant.
Its Hessian is therefore
\begin{equation}
 H_{\rm res}=\mathcal L-C^Tq_{\rm eff}C\geq0,\qquad
 H_{\rm res}D_f=0
 \label{dgb-residual-Hessian}
\end{equation}
for every fine infinitesimal gauge direction, with any
convenient fixed global base convention.
The nonnegativity uses the proved minimum comparison,
not stationarity alone.

Let $H$ be the actual derivative of the minimum section in
these relative coordinates. Since $CH=I$ and the residual
vanishes identically along that section, its positive
semidefinite Hessian obeys
\begin{equation}
                     H_{\rm res}H=0,\qquad
                     \mathcal LH=C^Tq_{\rm eff}.
 \label{dgb-residual-lift}
\end{equation}
No independently chosen flat lift is substituted for $H$.

A canonical completion is
\begin{equation}
                Q_b=H_{\rm res}+D_bD_b^T+C^TC>0.
 \label{dgb-canonical-completion}
\end{equation}
Indeed a null vector would satisfy $Cx=0$, $D_b^Tx=0$
and $x^T\mathcal Lx=0$. The constrained quotient positivity
then gives $x=D_b\alpha$, and $G_b\alpha=0$ implies $x=0$.
This is \eqref{dgb-completion} with $B=I-q_{\rm eff}$,
which need not be positive itself. No uniform lower
spectral bound for \eqref{dgb-canonical-completion} is
asserted as $K$ grows.

The general coarse compression can also be identified
exactly. Set
\begin{equation}
                   H_b=H-D_bG_b^{-1}D_b^TH .
 \label{dgb-adjusted-lift}
\end{equation}
Then $CH_b=I$, $D_b^TH_b=0$, $H_{\rm res}H_b=0$, and
$\mathcal LH_b=C^Tq_{\rm eff}$. For every completion
\eqref{dgb-completion}, it follows that
\begin{equation}
 Q_bH_b=C^T(q_{\rm eff}+B),\qquad
 q_{\rm eff}+B=H_b^TQ_bH_b>0.
 \label{dgb-effective-Schur}
\end{equation}
Multiplication by the inverses gives
\begin{equation}
 Q_b^{-1}C^T=H_b(q_{\rm eff}+B)^{-1},\qquad
 \boxed{\,M_b=(q_{\rm eff}+B)^{-1}\,}.
 \label{dgb-coarse-identification}
\end{equation}
Thus \eqref{dgb-main-identity} is equivalently
\begin{equation}
 D_{M,K}=\mathcal N_0\left[
      \log\det G_b-\frac12\log\det Q_b
                   +\frac12\log\det(q_{\rm eff}+B)\right].
 \label{dgb-effective-determinant}
\end{equation}
The coarse covariance is the inverse of the actual
relative-coordinate effective Hessian with its completion.
For the canonical choice $B=I-q_{\rm eff}$, it is exactly
$M_b=I$ and that coarse logarithm vanishes. The coarse
physics remains inside
$H_{\rm res}=\mathcal L-C^Tq_{\rm eff}C$ and hence in $Q_b$;
no uniform reference comparison is thereby obtained.

\subsection{A soft based-ghost mode}
\label{dgb-soft-ghost-section}

At the flat background $D_b$ is the nearest-neighbor gauge
gradient on scalar functions vanishing at the $K$-centers.
Take $K\geq3$ on a compatible unit-current torus and choose
one real color. Let $\varphi$ be zero at all centers and
one elsewhere. If there are $n_c$ centers, there are
$n_f=K^dn_c$ current sites, and exactly $2d$ incident
unit-difference bonds per center. Hence
\begin{equation}
 \|\varphi\|_2^2=(K^d-1)n_c,\qquad
 \|D_b\varphi\|_2^2=2dn_c,\qquad
 \lambda_{\min}(G_b)\leq\frac{2d}{K^d-1}.
 \label{dgb-Rayleigh}
\end{equation}
This is an upper Rayleigh bound, not a matching spectral
asymptotic. In four dimensions it implies
\[
                         \|G_b^{-1}\|\geq\frac{K^4-1}{8}.
\]
A final-physical-scale mass is $K^{-2}$ in the current
coordinates. Thus the estimate $\|G_b^{-1}\|\leq CK^2$
is false uniformly in $K$. After resetting final spacing
to one, $K^2G_b$ still has a Rayleigh eigenvalue tending
to zero.

This counterexample does not refute a bounded Maxwell
finite-part difference for the complete expression
\eqref{dgb-main-identity}. It rules out estimating the
based ghost separately as a uniformly massive operator.
The vector and coarse determinant must be kept coupled.

\subsection{An explicit full-site and coarse Schur decomposition}
\label{dgb-full-Schur-section}

Fix one final center as a global base. Let $D_f$ now denote
the fine gauge generator for all scalar sites except that
base, and split its domain into the non-$K$-centers and the
other $K$-centers:
\[
                            D_f=[\,D_b,D_e\,].
\]
These are disjoint site-coordinate columns, not orthogonal
images. The globally based action is free. Therefore
\begin{align}
 G_f&=D_f^TD_f>0,\notag\\
 S_g&=D_e^T(I-D_bG_b^{-1}D_b^T)D_e>0,\notag\\
 \det G_f&=\det G_b\,\det S_g .
 \label{dgb-ghost-Schur}
\end{align}
For any positive $Q_b$ as above, set
\begin{equation}
 Q_f=Q_b+D_eD_e^T,\qquad
 S_v=I+D_e^TQ_b^{-1}D_e.
 \label{dgb-vector-Schur-data}
\end{equation}
The determinant lemma gives
\begin{equation}
                         \det Q_f=\det Q_b\,\det S_v .
 \label{dgb-vector-Schur}
\end{equation}
Consequently the actual correction is exactly
\begin{equation}
 D_{M,K}=\mathcal N_0\left[
   \log\det G_f-\frac12\log\det Q_f
   -\log\det S_g+\frac12\log\det S_v
   -\frac12\log\det M_b\right].
 \label{dgb-full-decomposition}
\end{equation}
$S_g,S_v$ are coarse-site matrices; $M_b$ is a coarse-bond
matrix. The first two terms resemble a vector/ghost
combination, but all coarse corrections remain.
$G_f$ is globally based; it is not yet the translation-invariant
all-site massive ghost of a reference calculation.
If no nonbase coarse site remains, the corresponding
zero-dimensional determinants are one. For the canonical
completion $M_b=I$ by \eqref{dgb-coarse-identification};
the general formula retains all three coarse terms.

For \eqref{dgb-canonical-completion}, the vector operator is
\begin{equation}
 Q_f=s''+D_fD_f^T-\lambda\cdot f''
                         +C^T(I-q_{\rm eff})C.
 \label{dgb-vector-operator}
\end{equation}
If a proposed matched massive reference is
$P_\eta=s''+D_fD_f^T+\eta^2I$, with $\eta=K^{-1}$,
the exact difference is
\begin{equation}
 Q_f-P_\eta=-\lambda\cdot f''
                     +C^T(I-q_{\rm eff})C-\eta^2I.
 \label{dgb-reference-difference}
\end{equation}
The multiplier term vanishes at the flat background,
but its first and second background derivatives can
contribute to the determinant Hessian. Neither that
vanishing nor the coarse rank bounds its Maxwell
coefficient by itself. Equation
\eqref{dgb-reference-difference} is a comparison interface,
not an identification with an operator named in a
different background-field convention.
Any logarithmic comparison using $P_\eta$ additionally
requires its invertibility on the chosen finite germ.
No uniform such neighborhood is inferred here.

\subsection{A uniformly bounded coarse effective-Hessian factor}
\label{dgb-coarse-bound-section}

The general identification in
\eqref{dgb-coarse-identification} permits a uniform estimate
for one actual comparison term. Use the established strict
scalar family of the entire action $S_{MK}$, including its
pair-supported quadratic baseline. On a fixed small complex
ball it has uniform spatial bounds and an arbitrary
prescribed fixed cardinality reserve, independent of
$MK$ and compatible volume.

For $W=e^{iw}$, set
\[
                    \Xi_w(y)=\frac1i\log(e^{iy}e^{iw}).
\]
This is a per-bond analytic map with uniformly bounded
derivatives on a fixed small joint ball. With
$E_w=D_y\Xi_w(0)$ and $U_w=D_y^2\Xi_w(0)$,
\begin{equation}
 q_{\rm eff}(w)[u,v]
   =D^2S_{MK}(w)[E_wu,E_wv]
                 +DS_{MK}(w)[U_w(u,v)].
 \label{dgb-coarse-BCH}
\end{equation}
The background-current term is retained.
Compose each scalar support coefficient with $\Xi_w(y)$,
differentiate twice in $y$, and up to twice in $w$ using
fixed spare Cauchy radii. No support enlargement occurs.
Rows, columns and background source slots stay in the
same carrier. A fixed cardinality reserve absorbs their
finite counting factors.

Consequently, for some fixed $\mu,\rho>0$,
\begin{equation}
       \sup_{\|w\|_\infty<\rho}
                    \|q_{\rm eff}(w)\|_\mu\leq R_\mu,
 \label{dgb-coarse-Schur}
\end{equation}
where $\|\cdot\|_\mu$ is the maximum of the exponentially
weighted absolute block row and column sums.
The bound is uniform in $MK$ and volume. After reducing
the exponent to $m>0$, its source derivatives satisfy
\begin{align}
 \sum_{ij}e^{m[d(i,x)+d(j,x)]}
       |(\partial_xq_{\rm eff})_{ij}|&\leq C_1,\notag\\
 \sum_{ij}|(\partial_x\partial_yq_{\rm eff})_{ij}|
                              &\leq C_2e^{-md(x,y)}.
 \label{dgb-coarse-source-bounds}
\end{align}
To obtain these bounds, split the carrier-tree exponential
between its source separation and its row-to-source
distance, then sum that distance. The coarse cell fiber
has fixed dimension, with no factor from $K^4$ fine
coordinates.

\begin{proposition}[Bounded actual coarse determinant contribution]
\label{dgb-coarse-proposition}
Choose a fixed $a>2R_\mu$ and the completion $B=aI$.
The second-background-derivative kernel of
\[
                \tfrac12\mathcal N_0
                        \log\det(q_{\rm eff}+aI)
\]
has a uniform exponential row bound, independent of the
preceding depth, endpoint and compatible volume.
Its true infinite Maxwell coefficient, and its actual
finite reference-normalized coefficients above the
reference threshold, have a uniform absolute bound.
\end{proposition}

\begin{proof}
In the weighted Schur algebra, the Neumann series gives
\begin{equation}
 P_c=q_{\rm eff}+aI,\qquad
                         \|P_c^{-1}\|_\mu\leq 2/a.
 \label{dgb-coarse-inverse}
\end{equation}
The corresponding fine completion is positive on the
real germ: write $x=\xi+H_by$, $C\xi=0$.
Equation \eqref{dgb-effective-Schur} makes the cross term
zero, leaving the positive constrained-divergence form
in $\xi$ and $y^T(q_{\rm eff}+aI)y$.
No uniform fine spectral gap is inferred.

At finite volume the series
\begin{equation}
 \log\det P_c
   =(\dim P_c)\log a+
      \sum_{n\geq1}\frac{(-1)^{n+1}}n
                \operatorname{Tr}(a^{-1}q_{\rm eff})^n
 \label{dgb-coarse-log-series}
\end{equation}
converges absolutely. Let $K_c$ be the Hessian kernel of
$\mathcal N_0\log\det P_c$; the desired half-logarithm
has Hessian $K_c/2$. Explicitly,
\begin{equation}
 K_c(x,y)=\operatorname{Tr}\left[
       P_c^{-1}q_{{\rm eff},xy}
       -P_c^{-1}q_{{\rm eff},x}
                   P_c^{-1}q_{{\rm eff},y}\right].
 \label{dgb-coarse-trace}
\end{equation}
The first trace is bounded by the second line of
\eqref{dgb-coarse-source-bounds} and the inverse bound.
For the second trace, use the two source-rooted first
derivative bounds and spend a spare inverse-edge
exponent to connect their roots. Absolute summation gives
\begin{equation}
                            |K_c(x,y)|
                                \leq C_a e^{-m'd(x,y)}
 \label{dgb-coarse-Hessian-decay}
\end{equation}
for a fixed $m'>0$, uniformly in depth and volume.
Periodic sums are bounded by their infinite-lattice
counterparts. Its second displacement moment is therefore
uniformly bounded. The finite reference denominators
are uniformly bounded below above their established
threshold, which proves the finite statement as well.
\end{proof}

The actual action's local gauge identity makes its
relative Hessian transform by orthogonal bond-frame
conjugation under real local gauges. Thus the coarse
massive determinant has the appropriate total local
gauge identity on its common germ; no coefficientwise
Ward identity was used in the proof.
With $B=aI$, \eqref{dgb-effective-determinant} now reduces
the unbounded-endpoint comparison to the fine combination
\[
 \log\det G_b-\tfrac12\log\det\!
       \left[s''-\lambda\cdot f''
                     +D_bD_b^T+aC^TC\right],
\]
up to the proved uniformly bounded coarse contribution.
For a general completion the full-site vector difference
is $-\lambda\cdot f''+C^TBC-\eta^2I$. For this fixed-$a$
choice it is
\begin{equation}
              Q_f-P_\eta
                    =-\lambda\cdot f''+aC^TC-\eta^2I.
 \label{dgb-fixed-a-difference}
\end{equation}
Equation \eqref{dgb-reference-difference} instead records
the canonical completion $B=I-q_{\rm eff}$.
Each comparison uses all factors from its one chosen
completion.

\subsection{The remaining uniform comparison}
\label{dgb-scope-section}

The deposited outline asserts a bounded error in the
Maxwell coefficient when its composite determinant is
compared with a massive vector/ghost reference, at
$\eta=\ell^{-k}$. The exact identities above do not
establish that error bound. In particular, the formal
background effective-action convention and massless
nonzero-external-momentum amplitudes of a reference
calculation do not identify this constrained normalizer
automatically.

For a precisely matched reference, after removing the
bounded coarse effective-Hessian contribution, the
full-site formulation must still control the Maxwell
moment of the coarse-site package
\[
                         -\log\det S_g+\tfrac12\log\det S_v,
\]
together with the global-base and ghost-mass correction,
and the Maxwell moment of
$\log\det Q_f-\log\det P_\eta$ with all multiplier,
coarse-completion and mass terms for the selected
completion. They must also identify
the reference background with the actual smooth pre-gauge
minimum and preserve the actual finite extraction and
zero-mode conventions. Based-gauge covariance allows
evaluating \eqref{dgb-main-identity} at that pre-gauge
representative, but does not replace it by an arbitrary
external interpolation.

Each determinant Hessian has the concrete formula
\begin{equation}
 D^2\log\det P[u,v]
   =\operatorname{Tr}\left(
       P^{-1}P_{uv}-P^{-1}P_uP^{-1}P_v\right).
 \label{dgb-trace-interface}
\end{equation}
The established fixed-$K$ bounds make the relevant
second-displacement-moment trace sums meaningful. Their
constants may depend on $K$; they do not prove an $O(1)$
bound as $K\to\infty$.
The soft mode \eqref{dgb-Rayleigh} shows why a separate
uniform inverse estimate cannot supply this step.
A coupled cancellation or a scale-resolved comparison
is still required. No universal coefficient, sign,
averaged drift or ultraviolet coupling trajectory is
deduced here.
